\documentclass[11pt,a4paper]{article}

\usepackage[utf8]{inputenc}
\usepackage[T1]{fontenc}
\usepackage[french]{babel}
\usepackage{amsmath, amssymb, amsthm}
\usepackage{geometry}
\usepackage{booktabs}
\usepackage{listings}
\usepackage{xcolor}
\usepackage{graphicx}

\geometry{margin=2.5cm}

\title{Corrige - TD 3 : Analyse et Commande des systèmes linéaires}
\author{AU361 - Grenoble INP Esisar}
\date{Annee 2025-2026}

\begin{document}

\maketitle

\section*{Partie I - Analyse du processus à piloter}

Le système mécanique relie la commande $u$ à la vitesse $vit$ par un premier ordre de gain statique $G=5$ et de constante de temps $\tau=0.1$ s.

\subsection*{I.1. Expression de la fonction de transfert $H(p)$}

\textbf{Raisonnement :} La position $pos(t)$ est l'intégrale de la vitesse $vit(t)$. Dans le domaine de Laplace, on a $POS(p) = \frac{VIT(p)}{p}$. La fonction de transfert entre $U(p)$ et $VIT(p)$ est $\frac{5}{1+0.1p}$.

\textbf{Reponse :} Ainsi, la fonction de transfert du processus $H(p)$ est :
$$ 
\mathbf{H(p) = \frac{POS(p)}{U(p)} = \frac{5}{p(1+0.1p)}} 
$$

\subsection*{I.2. Valeur des pôles du processus}

\textbf{Raisonnement :} Les pôles annulent le dénominateur de $H(p)$ :
$$ 
p(1+0.1p) = 0 
$$

\textbf{Reponse :} 
$$ 
\mathbf{p_1 = 0 \quad \text{et} \quad p_2 = -10 \text{ rad/s}} 
$$

\subsection*{I.3. Stabilité asymptotique}

\textbf{Raisonnement :} Le système possède un pôle nul (à la limite de la stabilité), traduisant un comportement intégrateur pur.

\textbf{Reponse :} Le système n'est \textbf{pas asymptotiquement stable}.

\subsection*{I.4. Expressions temporelles des réponses indicielles}

\textbf{Raisonnement :} Pour un échelon de commande $U(p) = \frac{1}{p}$, on a :
\begin{itemize}
    \item Vitesse : $VIT(p) = \frac{5}{p(1+0.1p)} = 5 \left( \frac{1}{p} - \frac{0.1}{1+0.1p} \right) = \frac{5}{p} - \frac{5}{p+10}$.
    \item Position : $POS(p) = \frac{5}{p^2(1+0.1p)} = \frac{5}{p^2} - \frac{0.5}{p} + \frac{0.5}{p+10}$.
\end{itemize}

\textbf{Reponse :} En passant dans le domaine temporel :
\begin{itemize}
    \item Vitesse : $\mathbf{vit(t) = 5(1 - e^{-10t}) \cdot \Gamma(t)}$
    \item Position : $\mathbf{pos(t) = (5t - 0.5 + 0.5e^{-10t}) \cdot \Gamma(t)}$
\end{itemize}

\hrulefill

\section*{Partie II - Analyse du système asservi}

\subsection*{II.1. Fonction de transfert du régulateur}

\textbf{Reponse :} Pour un régulateur proportionnel de gain K :
$$ 
\mathbf{R(p) = K} 
$$

\subsection*{II.2. Fonction de transfert en boucle fermée (FTBF)}

\textbf{Raisonnement :} On applique la formule de la boucle fermée :
$$ 
FTBF(p) = \frac{R(p)H(p)}{1 + R(p)H(p)} = \frac{\frac{5K}{p(1+0.1p)}}{1 + \frac{5K}{p(1+0.1p)}} = \frac{5K}{0.1p^2 + p + 5K} 
$$
En mettant sous forme canonique du second ordre $\frac{1}{\frac{p^2}{\omega_n^2} + \frac{2\zeta}{\omega_n}p + 1}$ :
$$ 
FTBF(p) = \frac{1}{\frac{0.1}{5K}p^2 + \frac{1}{5K}p + 1} 
$$

\textbf{Reponse :} Par identification :
$$ 
\mathbf{\omega_n = \sqrt{50K} \quad \text{et} \quad \zeta = \frac{1}{\sqrt{2K}}} 
$$

\subsection*{II.3. Position en régime permanent (consigne échelon)}

\textbf{Raisonnement :} Sans perturbation ($\delta=0$) et pour une consigne échelon $cons = 1/p$, on applique le théorème de la valeur finale :
$$ 
pos(\infty) = \lim_{p \to 0} p \cdot FTBF(p) \frac{1}{p} = FTBF(0) = 1 
$$

\textbf{Reponse :} Le système suit parfaitement la consigne, \textbf{l'erreur de position est nulle}.

\subsection*{II.4. Erreur statique due à la perturbation}

\textbf{Raisonnement :} Avec $cons=0$ et $\delta=1/p$, le transfert entre la perturbation et la sortie est $T_\delta(p) = \frac{H(p)}{1+R(p)H(p)} = \frac{5}{0.1p^2 + p + 5K}$.
$$ 
pos(\infty) = \lim_{p \to 0} p \cdot T_\delta(p) \frac{1}{p} = T_\delta(0) = \frac{5}{5K} = \frac{1}{K} 
$$

\textbf{Reponse :} L'erreur statique due à cette perturbation est donc $\mathbf{\frac{1}{K}}$.

\subsection*{II.5. Minimisation de l'erreur statique sans dépassement}

\textbf{Raisonnement :} Pour minimiser l'erreur ($\frac{1}{K}$), il faut maximiser $K$. Cependant, pour garantir une réponse sans dépassement, il faut un amortissement $\zeta \ge 1$.
$$ 
\frac{1}{\sqrt{2K}} \ge 1 \implies 2K \le 1 \implies K \le 0.5 
$$
On choisit la valeur maximale $K = 0.5$.

\textbf{Reponse :} L'erreur statique due à la perturbation vaut alors :
$$
\mathbf{pos(\infty) = \frac{1}{0.5} = 2}
$$

\subsection*{II.6. Diagramme de Bode de $T_\delta(p)$ pour $K=2$}

\textbf{Raisonnement :} Pour $K=2$, $T_\delta(p) = \frac{5}{0.1p^2 + p + 10}$. Les paramètres deviennent $\omega_n = \sqrt{100} = 10$ rad/s et $\zeta = 0.5$. La pulsation de résonance est donnée par $\omega_r = \omega_n \sqrt{1 - 2\zeta^2}$.

\textbf{Reponse :} 
$$ 
\mathbf{\omega_r = 10 \sqrt{1 - 0.5} = \frac{10}{\sqrt{2}} \approx 7.07 \text{ rad/s}} 
$$

\subsection*{II.7. Réponse à une perturbation sinusoïdale}

\textbf{Raisonnement :} Fréquence $f = 50$ Hz $\implies \omega = 314$ rad/s. Amplitude $= 0.1$. En régime permanent, la position a l'allure d'une sinusoïde de même pulsation. L'amplitude $A$ de ces variations est $A = 0.1 \times |T_\delta(j\omega)|$.

\textbf{Reponse :} 
$$ 
\mathbf{A = \frac{0.5}{\sqrt{(5K - 0.1\omega^2)^2 + \omega^2}}} 
$$

\hrulefill

\section*{Partie III - Moteur asservi en position}

Le système intègre un correcteur $C(s)$, un transfert $\frac{K}{1+Ts}$ et un intégrateur $\frac{1}{s}$.

\subsection*{1. Expressions de $T_1(s)$ et $T_2(s)$}

\textbf{Raisonnement :} D'après le schéma, $V(s) = \frac{K}{1+Ts}[C(s)(CY(s) - Y(s)) - \Gamma_r(s)]$ et $Y(s) = \frac{V(s)}{s}$. En réorganisant :
$$ 
Y(s) \left[ 1 + \frac{K C(s)}{s(1+Ts)} \right] = \frac{K C(s)}{s(1+Ts)} CY(s) - \frac{K}{s(1+Ts)} \Gamma_r(s) 
$$

\textbf{Reponse :} On identifie les deux fonctions de transfert :
$$ 
\mathbf{T_1(s) = \frac{K C(s)}{s(1+Ts) + K C(s)} \quad \text{et} \quad T_2(s) = \frac{-K}{s(1+Ts) + K C(s)}} 
$$

\subsection*{2. Erreur en suivi de consigne échelon ($C(s) = K_p$)}

\textbf{Raisonnement :} Le transfert en boucle ouverte comporte un intégrateur pur ($\frac{1}{s}$). Le système est donc de classe 1.

\textbf{Reponse :} Ce qui garantit une erreur statique de position \textbf{nulle} pour une consigne échelon.

\subsection*{3. Erreur en réponse à une perturbation unitaire}

\textbf{Raisonnement :} L'erreur est $\epsilon = -y$ (puisque $cy=0$).
$$ 
\lim_{s \to 0} s \cdot T_2(s) \frac{1}{s} = T_2(0) = \frac{-K}{K K_p} = -\frac{1}{K_p} 
$$

\textbf{Reponse :} L'erreur statique est donc $\mathbf{\frac{1}{K_p}}$.

\subsection*{4. Erreur de traînage (consigne rampe)}

\textbf{Raisonnement :} Pour $CY(s) = 1/s^2$, l'erreur est $\epsilon(s) = CY(s)(1 - T_1(s)) = \frac{1+Ts}{s[s(1+Ts)+K K_p]}$.

\textbf{Reponse :} 
$$ 
\mathbf{\epsilon(\infty) = \lim_{s \to 0} s \epsilon(s) = \frac{1}{K K_p}} 
$$

\subsection*{5. Pas de surtension de gain ($K=4, T=0.25$)}

\textbf{Raisonnement :} $T_1(s)$ est un second ordre avec $\omega_n^2 = \frac{4K_p}{0.25} = 16K_p$ et $\frac{2\zeta}{\omega_n} = \frac{1}{4K_p}$, d'où $\zeta = \frac{1}{2\sqrt{K_p}}$. Il n'y a pas de résonance si $\zeta \ge \frac{1}{\sqrt{2}}$.

\textbf{Reponse :} 
$$ 
\mathbf{\frac{1}{2\sqrt{K_p}} \ge \frac{1}{\sqrt{2}} \implies \sqrt{K_p} \le \frac{\sqrt{2}}{2} \implies K_p \le 0.5} 
$$

\subsection*{6. Marge de phase de 45°}

\textbf{Raisonnement :} La FTBO est $L(j\omega) = \frac{4K_p}{j\omega(1+0.25j\omega)}$. Phase : $\phi = -90^\circ - \arctan(0.25\omega)$.
Pour $M_\phi = 45^\circ$, on a $\phi = -135^\circ \implies \arctan(0.25\omega) = 45^\circ \implies \omega = 4$ rad/s.
Le gain doit valoir 1 à cette pulsation :
$$ 
|L(j4)| = \frac{4K_p}{4\sqrt{1+1}} = \frac{K_p}{\sqrt{2}} = 1 
$$

\textbf{Reponse :} 
$$ 
\mathbf{K_p = \sqrt{2} \approx 1.41} 
$$
Puisque la phase ne dépasse jamais $-180^\circ$, la marge de gain est \textbf{infinie}.

\subsection*{7. Marge de module pour $K_p=1$}

\textbf{Raisonnement :} La marge de module $\Delta M$ est la distance minimale entre le lieu de Nyquist de $L(j\omega)$ et le point critique $-1$.
$$
\Delta M = \min_{\omega} |1 + L(j\omega)| = \min_{\omega} |S(j\omega)|^{-1}
$$

\textbf{Reponse :} En trouvant le minimum numérique, on obtient $\mathbf{\Delta M \approx 0.68}$.

\subsection*{8. Stabilité avec correcteur PI}

\textbf{Raisonnement :} Pour $C(s) = \frac{1+sT_a}{sT_b}$, l'équation caractéristique de la boucle fermée est :
$$ 
D(s) = T T_b s^3 + T_b s^2 + K T_a s + K = 0 
$$
Dressons le tableau de Routh :
\begin{table}[h!]
    \centering
    \begin{tabular}{ccc}
        \toprule
        \textbf{Puissance} & \textbf{Colonne 1} & \textbf{Colonne 2} \\
        \midrule
        $s^3$ & $T T_b$ & $K T_a$ \\
        $s^2$ & $T_b$ & $K$ \\
        $s^1$ & $K(T_a - T)$ & 0 \\
        $s^0$ & $K$ & 0 \\
        \bottomrule
    \end{tabular}
\end{table}
Pour que le système soit stable (en supposant $K>0$ et $T_b>0$), tous les éléments de la première colonne doivent être positifs.

\textbf{Reponse :} La condition de stabilité est donc $\mathbf{T_a > T}$.

\end{document}